\documentclass[11pt]{article}
\usepackage[margin=1in]{geometry}
\usepackage{amsmath,amssymb,amsfonts,mathtools}
\usepackage{array,booktabs,enumitem,graphicx,xcolor,hyperref}
\usepackage[T1]{fontenc}
\usepackage[utf8]{inputenc}
\title{The big picture}
\author{Source-backed arXiv sample}
\date{}
\begin{document}
\maketitle
\begin{abstract}
We discuss the conformal field theory and string field theory of the NSR superstring using a BRST operator with a nonminimal term, which allows all bosonic ghost modes to be paired into creation and annihilation operators. Vertex operators for the Neveu-Schwarz and Ramond sectors have the same ghost number, as do string fields. The kinetic and interaction terms are the same for Neveu-Schwarz as for Ramond string fields, so spacetime supersymmetry is closer to being manifest. The kinetic terms and supersymmetry don't mix levels, simplifying component analysis and gauge fixing.
\end{abstract}
\section{Source Metadata}
This sample is based on arXiv record \texttt{hep-th/9108021}. Authors: N. Berkovits, M.T. Hatsuda, and W. Siegel.
\section{Abstract Discussion}
We discuss the conformal field theory and string field theory of the NSR superstring using a BRST operator with a nonminimal term, which allows all bosonic ghost modes to be paired into creation and annihilation operators. Vertex operators for the Neveu-Schwarz and Ramond sectors have the same ghost number, as do string fields. The kinetic and interaction terms are the same for Neveu-Schwarz as for Ramond string fields, so spacetime supersymmetry is closer to being manifest. The kinetic terms and supersymmetry don't mix levels, simplifying component analysis and gauge fixing.
\end{document}
